Na camada de infraestrutura, cinco frentes ficam abertas. A primeira é completar o nó: substituir os dados sintéticos por sensores reais de temperatura e iluminância e colocar o microcontrolador em repouso entre transmissões, o que fecha os dois requisitos hoje parciais, RF1 e RNF1. A segunda é a medição de energia do nó por estado de operação, por transmissão e por cifragem, para a qual o firmware de bancada já está pronto; ela daria ao custo dos 16 bytes da \textit{tag} um valor medido em energia, e não apenas em tempo no ar. A terceira é o teste de alcance, com as placas afastadas, potência de transmissão maior e fatores de espalhamento de 7 a 12. É a condição em que o SNR sai da saturação e o RSSI passa a variar, ou seja, em que duas das variáveis dos modelos deixam de ser quase constantes. A quarta é a geração de tráfego anômalo real, em ambiente controlado, com um terceiro dispositivo que repita e forje quadros no canal do enlace. O \textit{gateway} já contabiliza em separado as rejeições por repetição e por autenticação, de modo que a coleta própria passaria a conter eventos rotulados e poderia ser usada na avaliação dos modelos de detecção, e não apenas como referência de operação normal. A quinta é a migração do \textit{backend} para infraestrutura em nuvem, com transporte cifrado entre \textit{gateway} e API: hoje os dois estão na mesma rede local e a requisição HTTP trafega em claro, autenticada apenas por um token.

Como extensão de prazo mais longo, uma implementação do Ascon em hardware reconfigurável (FPGA), usada como coprocessador do nó, permitiria repetir na mesma bancada a comparação com o acelerador de AES do microcontrolador, com as duas cifras em hardware.
